% concatenated from lewis_weights_manuscript_skeleton.tex, packages.tex, formatting_macros.tex, macros.tex
\documentclass[11pt]{article}

\usepackage[margin=1in]{geometry}
\usepackage{microtype}
\usepackage[T1]{fontenc}
\usepackage{lmodern}

\usepackage{amsmath,amssymb,amsthm,mathtools,xfrac}
\usepackage{bm}
\usepackage{mathrsfs}

\usepackage{algorithm}
\usepackage[noend]{algpseudocode}
\usepackage{booktabs}
\usepackage{graphicx}
\usepackage{caption}
\captionsetup[table]{font=small}
\captionsetup[figure]{font=small}

\usepackage[dvipsnames]{xcolor}

\colorlet{myblue}{RoyalBlue!70!Blue}
\colorlet{mymagenta}{Magenta!50!Plum}

\usepackage[
    colorlinks=true,
    linkcolor=myblue,
    citecolor=mymagenta,
    urlcolor=MidnightBlue
]{hyperref}

\usepackage[
    backend=biber,
    backref=true, 
    style=alphabetic,
    sorting=ynt,       
    maxbibnames=999,   
    maxcitenames=999, 
    maxalphanames=999, 
    giveninits=false   
]{biblatex}

\addbibresource{references.bib}

\usepackage{aliascnt}
\usepackage[nameinlink,capitalize,noabbrev]{cleveref}

\newcommand\numberthis{\addtocounter{equation}{1}\tag{\theequation}}

\renewcommand{\algorithmicrequire}{\textbf{Input:}}
\renewcommand{\algorithmicensure}{\textbf{Output:}}

\theoremstyle{plain}

\newtheorem{theorem}{Theorem}[section]

\newaliascnt{lemma}{theorem}
\newtheorem{lemma}[lemma]{Lemma}
\aliascntresetthe{lemma}

\newaliascnt{proposition}{theorem}
\newtheorem{proposition}[proposition]{Proposition}
\aliascntresetthe{proposition}

\newaliascnt{corollary}{theorem}
\newtheorem{corollary}[corollary]{Corollary}
\aliascntresetthe{corollary}

\newaliascnt{fact}{theorem}
\newtheorem{fact}[fact]{Fact}
\aliascntresetthe{fact}

\newaliascnt{claim}{theorem}
\newtheorem{claim}[claim]{Claim}
\aliascntresetthe{claim}

\newaliascnt{definition}{theorem}
\newtheorem{definition}[definition]{Definition}
\aliascntresetthe{definition}

\newaliascnt{assumption}{theorem}
\newtheorem{assumption}[assumption]{Assumption}
\aliascntresetthe{assumption}

\newaliascnt{remark}{theorem}
\newtheorem{remark}[remark]{Remark}
\aliascntresetthe{remark}

\numberwithin{equation}{section}


\crefname{equation}{}{}
\Crefname{equation}{}{}
\creflabelformat{equation}{#2(#1)#3}

\crefname{theorem}{Theorem}{Theorems}
\Crefname{theorem}{Theorem}{Theorems}

\crefname{lemma}{Lemma}{Lemmas}
\Crefname{lemma}{Lemma}{Lemmas}

\crefname{proposition}{Proposition}{Propositions}
\Crefname{proposition}{Proposition}{Propositions}

\crefname{corollary}{Corollary}{Corollaries}
\Crefname{corollary}{Corollary}{Corollaries}

\crefname{fact}{Fact}{Facts}
\Crefname{fact}{Fact}{Facts}

\crefname{claim}{Claim}{Claims}
\Crefname{claim}{Claim}{Claims}

\crefname{definition}{Definition}{Definitions}
\Crefname{definition}{Definition}{Definitions}

\crefname{assumption}{Assumption}{Assumptions}
\Crefname{assumption}{Assumption}{Assumptions}

\crefname{remark}{Remark}{Remarks}
\Crefname{remark}{Remark}{Remarks}

\crefname{algorithm}{Algorithm}{Algorithms}
\Crefname{algorithm}{Algorithm}{Algorithms}


\newenvironment{draftblock}{\begingroup\color{gray}}{\endgroup}

\renewcommand{\qedsymbol}{$\blacksquare$}
\usepackage{thm-restate}


\newcommand{\R}{\mathbb{R}}
\newcommand{\N}{\mathbb{N}}
\newcommand{\Sp}{\mathbb{S}_{+}}
\newcommand{\Spp}{\mathbb{S}_{++}}
\newcommand{\one}{\mathbf{1}}
\newcommand{\diag}{\operatorname{diag}}
\newcommand{\Diag}{\operatorname{Diag}}
\newcommand{\tr}{\operatorname{Tr}}
\newcommand{\Tr}{\operatorname{Tr}}
\newcommand{\rank}{\operatorname{rank}}
\newcommand{\range}{\operatorname{range}}
\newcommand{\supp}{\operatorname{supp}}
\newcommand{\eps}{\varepsilon}
\newcommand{\defeq}{\mathrel{\vcentcolon=}}
\newcommand{\eqdef}{\mathrel{=}\vcentcolon}
\newcommand{\ip}[2]{\left\langle #1,#2\right\rangle}
\newcommand{\norm}[1]{\left\lVert #1\right\rVert}
\newcommand{\abs}[1]{\left\lvert #1\right\rvert}
\newcommand{\paren}[1]{\left(#1\right)}
\newcommand{\brak}[1]{\left[#1\right]}
\newcommand{\Rmspos}{\mathbb{R}^m_{>0}}
\newcommand{\Rkpos}{\mathbb{R}_{>0}^k}
\newcommand{\Prob}{\operatorname{Prob}}
\newcommand{\E}{\mathbb{E}}
\newcommand{\jacobian}[2]{D_{#1}\left(#2\right)}
\newcommand{\dirderiv}[3]{D_{#1}\left(#2\right)[#3]}

\newcommand{\bA}{\mathbf{A}}
\newcommand{\bB}{\mathbf{B}}
\newcommand{\bBbar}{\overline{\mathbf{B}}}
\newcommand{\bC}{\mathbf{C}}
\newcommand{\bD}{\mathbf{D}}
\newcommand{\bE}{\mathbf{E}}
\newcommand{\bF}{\mathbf{F}}
\newcommand{\bG}{\mathbf{G}}
\newcommand{\bH}{\mathbf{H}}
\newcommand{\bI}{\mathbf{I}}
\newcommand{\bJ}{\mathbf{J}}
\newcommand{\bK}{\mathbf{K}}
\newcommand{\bKbar}{\overline{\mathbf{K}}}
\newcommand{\Kbar}{\overline{K}}
\newcommand{\bM}{\mathbf{M}}
\newcommand{\bP}{\mathbf{P}}
\newcommand{\bPbar}{\overline{\mathbf{P}}}

\newcommand{\bq}{\mathbf{q}}
\newcommand{\bQ}{\mathbf{Q}}
\newcommand{\bW}{\mathbf{W}}
\newcommand{\bV}{\mathbf{V}}
\newcommand{\bX}{\mathbf{X}}
\newcommand{\bY}{\mathbf{Y}}
\newcommand{\bzero}{\mathbf{0}}

\newcommand{\ba}{\mathbf{a}}
\newcommand{\bb}{\mathbf{b}}
\newcommand{\bh}{\mathbf{h}}
\newcommand{\bw}{\mathbf{w}}
\newcommand{\bv}{\mathbf{v}}
\newcommand{\bx}{\mathbf{x}}
\newcommand{\by}{\mathbf{y}}
\newcommand{\bu}{\mathbf{u}}
\newcommand{\bxbar}{\overline{\mathbf{x}}}
\newcommand{\bwhat}{\widehat{\mathbf{w}}}
\newcommand{\bwstar}{\mathbf{w}^\star}
\newcommand{\bwbar}{\overline{\mathbf{w}}}
\newcommand{\bWstar}{\mathbf{W}^\star}
\newcommand{\bMstar}{\mathbf{M}^\star}
\newcommand{\bMbar}{\overline{\mathbf{M}}}
\newcommand{\bAblk}[1]{\mathbf{A}_{[#1]}}
\newcommand{\bGbar}{\overline{\mathbf{G}}}
\newcommand{\bBblk}[1]{\mathbf{B}_{[#1]}}
\newcommand{\bSigma}{\bm{\Sigma}}
\newcommand{\bpi}{\bm{\pi}}

\newcommand{\calL}{\mathcal{L}}
\newcommand{\wstar}{w^{\star}}
\newcommand{\what}{\widehat{w}}
\newcommand{\wbar}{\overline{w}}
\newcommand{\Wbar}{\overline{\mathbf{W}}}
\newcommand{\bWbar}{\overline{\mathbf{W}}}
\newcommand{\twosqdel}{2\sqrt{\delta}}

\newcommand{\diagex}{\mathbf{D}(\mathbf{x})} 
\newcommand{\projex}{\mathbf{P}(\mathbf{x})}

\NewDocumentCommand{\sens}{g}{%
  s^{(p)}\IfValueT{#1}{_{#1}}%
}
\newcommand{\maxsens}{s_{\max}^{(p)}}

\NewDocumentCommand{\levsc}{g}{%
  \ensuremath{\sigma\IfValueT{#1}{_{#1}}}%
}

\NewDocumentCommand{\blevsc}{g}{%
  \ensuremath{\sigma\IfValueT{#1}{_{[#1]}}}%
}

\newcommand{\calT}{\mathcal{T}}
\newcommand{\omrp}{\gamma} 
\newcommand{\kldiv}[2]{\mathcal{D}_{\mathrm{KL}}\left(#1 \,\|\, #2\right)}
\newcommand{\KL}[2]{D_{\mathrm{KL}}\left(#1 \,\|\, #2\right)} 
\newcommand{\logdetdiv}[2]{\mathcal{D}_{-\log\det}\left(#1, #2\right)}


\newcommand{\todo}[1]{\textcolor{gray}{\textbf{[TODO:} #1\textbf{]}}}
\newcommand{\pswt}[1]{\textcolor{magenta}{\textbf{[SP:} #1\textbf{]}}}


\title{Computing Lewis Weights to High Precision by Fixed-Point Iteration}
\author{Swati Padmanabhan
\thanks{Department of Industrial and Systems Engineering, University of Minnesota, Twin Cities. pswt@umn.edu}}
\date{}

\begin{document}

\maketitle


\begin{abstract}
The $\ell_p$-Lewis weights of a matrix are defined by a fixed-point equation. For $p<4$, Cohen and Peng~\cite{cohen2015lp} showed that iterating an equivalent rearrangement of this equation computes Lewis weights to high precision; for $p\geq4$, prior high-precision methods instead use optimization-based approaches. We show that the direct Lewis fixed-point iteration, appearing in~\cite{lee2016faster}, computes Lewis weights to high precision for every $p>2$.

For a matrix $\bA\in\mathbb R^{m\times n}$ partitioned into row blocks $\bA_{[1]},\ldots,\bA_{[k]}$, we compute, for $p>2$, coordinatewise $\varepsilon$-approximate $\ell_p$ block Lewis weights in $O\left(p\log\frac{p\sqrt{ \sum_{i=1}^k\rank(\bA_{[i]}) }}{\varepsilon}\right)$ 
rounds of exact leverage-score-vector computations. For ordinary Lewis weights, this becomes
$O\left(p\log\frac{p\sqrt m}{\varepsilon}\right)$, improving the $O(p^2\log(m/\varepsilon))$ bound of Gribling, Sidford, and Zhang~\cite{pmlr-v336-gribling26a} for $p\geq4$.

\looseness=-1Our main observation is that each direct Lewis update contracts the KL divergence to the true weights by a factor of $1-\frac{2}{p}$. We also give an alternate explanation of this contraction through volume sampling and entropic independence. Synthetic experiments closely match our predicted local contraction rates and iteration counts grow approximately linearly with $p$; real-data experiments illustrate the information-concentration tradeoff of finite-$p$ block Lewis designs.
\end{abstract}
\newpage 
\section{Introduction}
\label{sec:intro}

The $\ell_p$ Lewis weights of a matrix are defined via a fixed-point equation \Cref{eq:Lewis-fixed-point}; this raises a natural computational question: can the corresponding fixed-point iteration be directly used to compute Lewis weights to high precision? 
The work of~\cite{cohen2015lp} answers this question affirmatively for $p<4$ by iterating an equivalent rearrangement of \Cref{eq:Lewis-fixed-point}. 
For $p\geq4$, however, the analysis of this method in~\cite{cohen2015lp} breaks down, and subsequent high-precision algorithms instead rely on optimization-based approaches~\cite{cohen2015lp,lee2016faster,lee2019solving,fazel2022computing,pmlr-v336-gribling26a}.
We show that the fixed-point iteration directly associated with \Cref{eq:Lewis-fixed-point}, appearing in~\cite{lee2016faster}, can, in fact, be used to compute Lewis weights to high precision for $p>2$, thereby answering this question affirmatively in this regime as well.

\begin{definition}[Lewis weights~\cite{lewis1978finite,cohen2015lp}]\label{def:LW}
Given a full-rank matrix $\bA\in\R^{m\times n}$ with $m\geq n$ and no zero rows, and a scalar $p\in(0,\infty)$, the $\ell_p$-Lewis weights of $\bA$ are the entries of the unique\footnote{Existence and uniqueness were first proven by D.R. Lewis~\cite{lewis1978finite}} vector $\wbar\in\Rmspos$ satisfying the equation
\[
\wbar_i=\levsc{i}\left(\Wbar^{1/2-1/p}\bA\right)
\text{ for all }i\in[m], \numberthis\label{eq:Lewis-fixed-point}
\]
where $\ba_i^\top$ is the $i^\mathrm{th}$ row of $\bA$, $\Wbar$ the diagonal matrix with vector $\bwbar$ along its diagonal, and $\levsc{i}(\bB):=\bb_i^\top(\bB^\top\bB)^{-1}\bb_i$ denotes the $i^\mathrm{th}$ leverage score of a full-column-rank matrix $\bB$. 
\end{definition}

\looseness=-1Lewis weights have found applications in $\ell_p$ sampling and regression~\cite{cohen2015lp,musco2022active} and in the construction of self-concordant barriers and related sampling algorithms~\cite{lee2019solving,laddha2020strong}.
Beyond the (row-wise) definition above, this fixed-point construction has a natural block analogue, in which a single weight is assigned to each block according to its total leverage after blockwise rescaling.
In this paper, we restrict ourselves to $p>2$, which is therefore also the regime for which we state \Cref{def:blockLW}.
\begin{definition}[Block Lewis weights~\cite{jambulapati2023sparsifying}]\label{def:blockLW} 
Given a full rank matrix $\bA\in\R^{m \times n}$ whose rows are partitioned into $k$ nonzero blocks $\bA_1, \bA_2, \dotsc, \bA_k$, and a scalar $p >2$, the block $\ell_p$-Lewis weights of $\bA$ are the entries of the  unique\footnote{Existence follows from the block-Lewis construction of~\cite{jambulapati2023sparsifying}; uniqueness for $p>2$ also follows from \Cref{lem:kl-contraction}.}
 vector $\wbar\in\Rkpos$ satisfying \[\wbar_i = \blevsc{i}\left(\Wbar^{1/2-1/p}\bA\right) \text{ for all } i \in [k],\] where $\bWbar:=\Diag(\wbar_1\bI_{m_1},\dotsc,\wbar_k\bI_{m_k})$, and
$
\blevsc{i}(\bB)
:=
\sum_{j \in I_i}\levsc{j}(\bB)
$
denotes the total leverage score of the $i^\mathrm{th}$ block of $\bB$.
\end{definition} 

When each block in \Cref{def:blockLW} comprises just one row, \Cref{def:blockLW} collapses to the (row-wise) Lewis weights from \Cref{def:LW}. 
Block Lewis weights were introduced in~\cite{jambulapati2023sparsifying} for sparsifying sums of norms, and have since been further developed for matrix block-norm approximation and distributionally robust regression~\cite{manoj2025change,manoj2026distributionally}.
We use the block formulation as a common setting in which to study the fixed-point dynamics of both row-wise and block Lewis weights.

\subsection{Our contribution}
\label{sec:intro-contribution}
In this subsection, we give an overview of our theoretical and experimental results. 

\paragraph{(1) Theoretical results.}  
Our main result is a guarantee for computing $\varepsilon$-approximate block Lewis weights for every $p>2$. 
Our algorithm (displayed in \Cref{alg:main}) iterates the fixed-point equation in \Cref{def:blockLW}; its row-wise version appears in~\cite[Algorithm~6]{lee2016faster}, where averaging the iterates yields one-sided approximate Lewis weights.  
We show that, with appropriate initialization and final output conversion, this iteration yields high-precision block Lewis weights.

\begin{theorem}[Block Lewis weights]\label{thm:blkLW}
Let $\bA\in\R^{m\times n}$ be a full-column-rank matrix partitioned into $k$ nonzero blocks $\bAblk{1}, \bAblk{2}, \ldots,\bAblk{k}$ with $\bAblk{i}\in\R^{m_i\times n}$ for each $i\in [k]$. Assume  $p>2$, and let $\bwbar\in\Rkpos$ denote the $\ell_p$-block-Lewis weights of $\bA$. Given an accuracy parameter $\eps\in(0,1)$, there is an algorithm (\Cref{alg:main}) which returns $\bwhat\in\Rkpos$ satisfying
\[
  (1-\eps)\wbar_i
  \le \what_i
  \le (1+\eps)\wbar_i
  \qquad\text{for every }i\in[k],
\]
in $O\left(
    p\log\left(\frac{p \sqrt{\sum_{i=1}^k \rank\left(\bAblk{i}\right)}}{\eps}\right)\right)$  iterations. Each iteration computes all $m$ exact row leverage scores of a matrix $\bD\bA$, where $\bD = \Diag\left(d_1\bI_{m_1}, d_2\bI_{m_2}, \dotsc, d_k\bI_{m_k}\right)$ for some positive $d_1, d_2, \dotsc, d_k$. 
\end{theorem}

\noindent Specializing each block to a single row gives the following guarantee. 
\begin{corollary}
\label{cor:ordinary-main}
Let  $\bA\in\R^{m\times n}$ be a full-column-rank matrix with no zero rows, let $p>2$, and let $\bwbar\in\Rmspos$ be the vector of its $\ell_p$-Lewis weights. Given an accuracy parameter $\eps\in(0,1)$, there is an algorithm which returns $\bwhat\in\Rmspos$ satisfying
\[
  (1-\eps)\wbar_i
  \le \what_i
  \le (1+\eps)\wbar_i \quad
\text{ for every }i\in[m],
\]
in $O\left(
    p\log\left(\frac{p\sqrt{m}}{\eps}\right)
  \right)$ iterations. Each iteration computes all $m$ exact row leverage scores of $\bD\bA\in\R^{m\times n}$ where $\bD = \Diag\left(d_1, d_2, \dotsc, d_m\right)$ for some positive $d_1, d_2, \dotsc, d_m$. 
\end{corollary}
For $p\geq4$, \Cref{cor:ordinary-main} gives, to our knowledge, the first nearly dimension-free high-precision guarantee for the direct Lewis fixed-point iteration. 
In earlier work, Cohen and Peng~\cite{cohen2015lp} show that an equivalent rearrangement of the Lewis fixed-point equation is contractive in coordinatewise multiplicative error only for $p<4$ (see \Cref{sec:intro-approachOverview}). Our key observation is that the direct Lewis update nevertheless has a global \emph{KL contraction}: each update contracts the KL divergence to the true Lewis weights by a factor of $1-\frac{2}{p}$, which in turn drives our $\widetilde{O}(p)$ rate.

In addition to a direct proof of this contraction in \Cref{sec:kl-decay}, we  present in \Cref{sec:global-intuition} an alternate proof through volume sampling and entropic independence~\cite{derezinski2021determinantal,anari2022entropic}. We believe this  alternate perspective  makes the origin of the contraction factor $1-\frac{2}{p}$ more intuitive.

While the above contraction yields the stated \emph{global} iteration bounds, it does not characterize the asymptotic convergence rate near the fixed point $\bwbar$.
We therefore also study the \emph{local} dynamics of the row-wise iteration by linearizing the Lewis update at $\bwbar$.
The spectral radius of the resulting Jacobian predicts the asymptotic local contraction rate, which we  corroborate experimentally. 

\begin{remark} For $p \geq4$, \Cref{cor:ordinary-main} improves upon the $\widetilde{O}(p^2)$ rate of \cite{pmlr-v336-gribling26a}. Our total runtime is $O\left(p mn^{\omega-1}\log\left(\frac{p\sqrt{m}}{\varepsilon}\right)\right)$, where $\omega$ denotes the matrix-multiplication exponent~\cite{huang1998fast, alman2025more}. Each iteration computes one exact leverage-score vector in $O(mn^{\omega-1})$ operations. Our analysis is in exact real arithmetic, and we make no claims about bit complexity, approximate leverage scores, or floating-point stability. Our result has implications in the estimation of $\ell_p$ sensitivities, regularized D-optimal design, and the evaluation of the Lewis weights barrier, all of which we describe in \Cref{sec:apps}.
\end{remark} 

\paragraph{(2) Experimental results.}
We complement our theoretical results with synthetic and real-data experiments. 
Our synthetic experiments first test the local fixed-point dynamics mentioned above. 
Near the fixed point, the observed contraction rates closely agree with the rates predicted by our local analysis (\Cref{fig:local-dyn}). 
At the level of overall iteration counts, the number of iterations grows approximately linearly with $p$ for row-wise Lewis weights on Gaussian matrices (\Cref{fig:A1}), and this qualitative behaviour persists on a structured, leverage-heterogeneous matrix family (\Cref{fig:structured-ordinary}). 
We also test the fixed-point iteration in floating-point arithmetic on ill-conditioned matrices with identical exact Lewis weight trajectories; the observed iteration counts remain essentially unchanged over a broad range of conditioning, until finite precision imposes an accuracy floor (\Cref{fig:a4-map-error,fig:a4-residual-floor}). 
The block Lewis weight iteration exhibits an analogous dependence on $p$ (\Cref{fig:block-scaling}). 

The simplicity of our algorithm also makes it easy to compute block Lewis weights across a range of values of $p$ on real datasets. 
Motivated by the connection between Lewis weights and regularized D-optimal design (cf. \Cref{sec:apps}), we study these designs on the Intel Berkeley Research Lab sensor dataset~\cite{guestrin2004distributed} and meteorological data from the NOAA National Data Buoy Center. 
As $p$ increases, the finite-$p$ designs interpolate from diffuse, near-uniform allocations toward the more concentrated D-optimal design, and can recover much of the available D-optimal information gain before becoming comparably concentrated (\Cref{fig:design-geometry}). 
We then test these designs under different patterns of block unavailability. 
When few blocks are unavailable, larger values of $p$ preserve more information; as the number of unavailable blocks increases, the values of $p$ preserving the most information shift systematically downward (\Cref{fig:block-failure}).
Thus, our experiments illustrate both the dynamics of the fixed-point iteration and the tradeoff between information and concentration exhibited by the finite-$p$ block Lewis designs it computes.

\subsection{Prior work}\label{sec:intro-approachOverview}
Prior work on computing Lewis weights to high precision exhibits a sharp difference between the regimes $p<4$ and  $p\geq 4$.

\paragraph{The regime $p<4$.}
For $p<4$, Cohen and Peng~\cite{cohen2015lp} analyze an equivalent rearrangement of \Cref{eq:Lewis-fixed-point}. 
In particular, one update of their algorithm takes the form
\[
w_i^+
:=
\left(
\ba_i^\top
\left(\bA^\top\bW^{1-2/p}\bA\right)^{-1}
\ba_i
\right)^{p/2}.
\]
They show that if the current weights are within a multiplicative factor $c$ of the true Lewis weights, then one such update improves this factor to $c^{|p/2-1|}$. 
Thus, the coordinatewise multiplicative error contracts when $|\frac{p}{2}-1|<1$, or equivalently when $p<4$; at $p=4$ the factor reaches $1$, and for $p>4$ it exceeds $1$. 
Consequently, for $p\geq4$, this analysis does not prove high-precision convergence.

\paragraph{The regime $p\geq4$.} 
For $p\geq4$, subsequent algorithms instead turned to optimization formulations of Lewis weights. 
~\cite{cohen2015lp} use the determinant-maximization formulation together with cutting-plane methods; Lee~\cite{lee2016faster} develops a mirror-descent-based approach; and Lee and Sidford~\cite{lee2019solving} optimize a volumetric potential, using a homotopy argument to obtain a sufficiently good initialization. 
These methods establish efficient computation in this regime, but do not give a nearly dimension-free number of leverage-score computations with logarithmic dependence on $\frac{1}{\eps}$.

\paragraph{High precision for $p\geq4$.} 
Fazel, Lee, Padmanabhan, and Sidford~\cite{fazel2022computing} gave the first algorithm for $p\geq4$ that computes Lewis weights to high precision using a nearly dimension-free number of leverage-score-vector computations, using
$O\left(p^3\log\left(\frac{mp}{\eps}\right)\right)$
such computations. 
Their algorithm minimizes a regularized log-determinant objective via a quasi-Newton-type update. 
Sufficient progress is guaranteed under a certain rounding condition, which is either restored through a separate rounding procedure or maintained by choosing coordinate-dependent step sizes. 
Gribling, Sidford, and Zhang~\cite{pmlr-v336-gribling26a} subsequently improved the dependence on $p$ from cubic to quadratic, obtaining
$O\left(p^2\log\left(\frac{m}{\eps}\right)\right)$
leverage-score-vector computations. 
Their analysis uses local relative smoothness to analyze gradient-based methods for convex formulations of Lewis weights.

\paragraph{Block Lewis weights.}
Block Lewis weights were introduced by Jambulapati, Lee, Liu, and Sidford~\cite{jambulapati2023sparsifying} in the study of sparsifying sums of norms and further developed by Manoj and Ovsiankin~\cite{manoj2025change} for matrix block-norm sparsification and change-of-measure arguments. 
The algorithms in these works compute sampling weights or block Lewis overestimates sufficient for these applications, rather than coordinatewise high-precision approximations to the block Lewis weights, which is our focus in \Cref{thm:blkLW}.

\paragraph{Our approach.}
Departing from the convex optimization formulations used in the high-precision algorithms above for $p\geq4$, we return to the fixed-point iteration directly associated with \Cref{eq:Lewis-fixed-point}. 
With a suitable initialization and final output conversion, our algorithm (\Cref{alg:main}) repeatedly applies this fixed-point update. 
We show that, for $p>2$, each update contracts the KL divergence to the true Lewis weights, thereby allowing the direct fixed-point approach to extend beyond $p=4$.

\subsection{Applications}\label{sec:apps}

Lewis weights play several conceptually distinct roles in algorithms and optimization, from data-dependent importance scores, to experiment design, to self-concordant barrier construction. Below, we review a small subset of these connections. 

 
 
\paragraph{Sensitivity estimation.}
Sensitivity sampling is an importance-sampling framework in which data points are sampled according to the extent to which they can influence the objective, and has been used extensively in  coresets~\cite{langberg2010universal, varadarajan2012sensitivity}, clustering~\cite{feldman2011unified}, logistic regression~\cite{munteanu2018coresets}, and $\ell_p$ subspace embeddings~\cite{wy2023ICML}. 
The $\ell_p$ sensitivity of the $i^\mathrm{th}$ row of $\bA$ is defined as
$\sens{i}(\bA) := \sup_{\bx\neq \bzero} \frac{|\ba_i^\top \bx|^p}{\|\bA\bx\|_p^p}$.  
Unlike Lewis weights, computing sensitivities directly generally requires solving a separate maximization problem for each row, and recent work has therefore studied efficient methods for approximating them~\cite{padmanabhan2023computing}. 
 Lewis weights are closely related to sensitivities and can therefore serve as efficiently computable surrogates for them. In particular, for $p>2$, it is known that $\sens{i}(\bA) \leq n^{p/2-1}\wbar_i$ \cite[Fact~3.3]{padmanabhan2023computing}. Via a simple ``witness'' argument similar in spirit to that in
\cite[Fact~A.1]{padmanabhan2023computing}, we additionally observe that
the reverse inequality holds as well:
 
\begin{restatable}{proposition}{sensprop}\label{prop:sens}
Let $\bA\in\R^{m\times n}$ have full column rank and no zero rows, let $p>2$, and let $\bwbar\in\Rmspos$ denote its $\ell_p$-Lewis weights. Then, for every $i\in[m]$, we have  $\wbar_i\leq \sens{i}(\bA)$. Combined with the known bound $\sens{i}(\bA)\leq n^{p/2-1}\wbar_i$, this gives
\[\wbar_i\leq \sens{i}(\bA)\leq n^{p/2-1}\wbar_i.\]
Consequently, if $\bwhat$ is an $\eps$-approximation to $\bwbar$, i.e.,
$(1-\eps)\wbar_i\leq\what_i\leq(1+\eps)\wbar_i$ for every $i\in[m]$, then, simultaneously for every $i\in[m]$, we have 
$\frac{\what_i}{1+\eps}
\leq
\sens{i}(\bA)
\leq
\frac{n^{p/2-1}\what_i}{1-\eps}.$
\end{restatable}
 
 The first inequality in \Cref{prop:sens} turns the previously known one-sided Lewis weight surrogate for sensitivities into a two-sided comparison. Combined with \Cref{cor:ordinary-main}, it yields simultaneous bounds on all $m$ sensitivities using only a high-precision Lewis weight computation, with no additional individual sensitivity evaluations. In particular, defining $\maxsens(\bA):=\max_{i\in[m]}\sens{i}(\bA)$ and taking maxima in \Cref{prop:sens} gives an $n^{p/2-1}(1+O(\eps))$-factor approximation to $\maxsens(\bA)$. The best prior general-purpose guarantee we are aware of for estimating maximum $\ell_p$ sensitivity without explicitly computing every individual sensitivity is the $O(n^{p/2})$ factor of~\cite{padmanabhan2023computing} for $p>2$. Thus, our guarantee improves the dimension dependence of the approximation factor. 

\paragraph{Regularized D-optimal design.}
D-optimal experimental design is a classical problem in statistics  in which resources are allocated among a collection of experiments so as to maximize the determinant of the resulting information matrix ~\cite{pukelsheim2006optimal,pmlr-v70-allen-zhu17e,pmlr-v99-madan19a,singh2020approximation}. 
It is closely connected to the John ellipsoid problem~\cite{john1948extremum}, which is dual to a relaxed D-optimal-design problem~\cite{Todd2016}. 

For $p>2$, computing the $\ell_p$-Lewis weights is equivalent, after a change of variables, to solving a $\frac{p}{p-2}$-regularized D-optimal-design problem~\cite{cohen2015lp,lee2019solving,fazel2022computing}. 
The regularizer incentivizes spreading the allocation across experiments ---  as $p$ increases, the exponent $\frac{p}{p-2}$ decreases toward $1$, weakening the incentive, and the solution approaches the unregularized D-optimal design. 
This  interpolation between concentrated and uniform allocations motivates our experiments in \Cref{sec:real-data-expts}, where we empirically study how finite-$p$ block Lewis designs vary with $p$.

In light of the above connection, \Cref{cor:ordinary-main} gives a high-precision solver for this regularized design problem using $
O\left(p\log\frac{p\sqrt m}{\delta}\right)$ 
exact leverage-score-vector computations, where $\delta$ denotes the desired coordinatewise accuracy of the corresponding Lewis weights. 

\paragraph{Lewis barriers and Dikin walks.}
Self-concordant barriers are a foundational tool in convex optimization~\cite{nesterov1994interior} and also underlie many algorithms for sampling from polytopes. 
For polytopes, \cite{lee2019solving} introduced a nearly optimal self-concordant barrier based on $\ell_p$-Lewis weights with $p=\Theta(\log m)$; in particular, their barrier is $O(n\log^5 m)$-self-concordant. Lewis weight barriers and Dikin-type walks have since been used in many sampling algorithms~\cite{laddha2020strong,gatmiry2023sampling,jiang2024regularized,kook2024gaussian}. For the~\cite{lee2019solving} barrier, accurately computing the underlying Lewis weights is an important subroutine. Specializing \Cref{cor:ordinary-main} to $p=\Theta(\log m)$ reduces this computation to 
$O\left(\log m\log\frac{\sqrt m\log m}{\delta}\right)$
exact leverage-score-vector computations for coordinatewise relative accuracy $\delta$. 

\subsection{Notation and preliminaries}\label{sec:intro-notationPrelims}

Throughout, we assume $p>2$. We use boldface to denote non-scalar quantities, with boldfaced lowercase letters for vectors and boldfaced uppercase letters for matrices. We use $\bA$ to denote our full-column-rank $m\times n$ ($m\geq n$) real-valued input matrix. When discussing block Lewis weights, we partition the rows of $\bA$ into $k$ nonzero blocks, using $\bAblk{i}\in\R^{m_i\times n}$ for the $i^\mathrm{th}$ block. We use $\bwbar$ for the vector of Lewis weights of $\bA$; these weights are either row-wise or block-wise, as will be clear from context: $\bwbar\in\Rmspos$ are the row-wise weights and $\bwbar\in\Rkpos$ the block-wise weights.

We use $\levsc{i}(\bB)$ to denote the $i^\mathrm{th}$ leverage score of a full-column-rank matrix $\bB$. When its rows are partitioned into blocks as above, we denote the $i^\mathrm{th}$ block leverage score by $\blevsc{i}(\bB):=\Tr\left(\bBblk{i}(\bB^\top\bB)^{-1}\bBblk{i}^\top\right)$. We have a slight notational overloading for diagonal matrices: for a row-weight vector (e.g., $\bw\in\R^m$), we use its uppercase boldfaced form ($\bW$ in this example) to denote the diagonal matrix $\bW_{ii}=w_i$; in the block setting, for a block-weight vector $\bw\in\R^k$, we instead use $\bW=\Diag\left(w_1\bI_{m_1},w_2\bI_{m_2},\dotsc,w_k\bI_{m_k}\right)$, so that the weight associated with a block is repeated on each of its rows. The intended meaning will be clear from the dimension of the weight vector. We use $w_i^{(t)}$ to denote the $i^\mathrm{th}$ entry of $\bw^{(t)}$, where the superscript $t$ indexes the algorithm's iteration.

We use $\one$ to denote the vector of all ones, $\bI$ for the identity matrix, and $\bzero$ for the zero vector or matrix;  sizes are implied by context. For compatible symmetric matrices $\bX,\bY$, we use $\bX\succeq\bY$ to mean that $\bX-\bY$ is positive semidefinite. For a positive integer $s$, we use $[s]$ for the set $\{1,2,\dotsc,s\}$. For vectors, the operations of exponentials, logarithms, and powers are applied coordinatewise.


The identities in \Cref{fact:block-levsc} are block analogues of classical facts about leverage scores, which we provide for completeness (with a proof in \Cref{sec:app-technical-lemmas}). In particular, the second identity specializes, for graph incidence matrices, to Foster's theorem~\cite{foster1949average}. The formulation of block leverage scores we use previously appeared in~\cite{KLPSS16,xu2016sub} and has been widely used in several forms in, e.g., matrix sparsification and numerical linear algebra.

\begin{restatable}{fact}{blocklevscfact}\label{fact:block-levsc}
Let $\bB:=\begin{bmatrix}\bB_{[1]}^\top&\bB_{[2]}^\top&\cdots&\bB_{[k]}^\top\end{bmatrix}^\top\in\R^{m\times n}$ have full column rank, with $\bB_{[i]}\in\R^{m_i\times n}$, and define $\blevsc{i}:=\Tr\left(\bB_{[i]}(\bB^\top\bB)^{-1}\bB_{[i]}^\top\right)$. Then:
\begin{enumerate}
\item $0\leq\blevsc{i}(\bB)\leq\rank(\bB_{[i]})$ for every $i\in[k]$.
\item $\sum_{i=1}^k\blevsc{i}(\bB)=n$.
\end{enumerate}
\end{restatable}

\section{Algorithm and analysis}
\label{sec:main-results}
In this section, we describe our algorithm and analysis that guarantee \Cref{thm:blkLW}. Our algorithm (displayed in \Cref{alg:main}) is a slight modification of \cite[Algorithm 6]{lee2016faster}. Given the $k$-block-partitioned matrix $\bA\in\R^{m\times n}$, the parameter $p>2$, and an accuracy $\eps\in(0,1)$, it first computes $r_i:=\rank(\bAblk{i})$ and $R:=\sum_{i=1}^k r_i$. When $R=n$, the  exact weights are $\wbar_i=r_i$, which the algorithm therefore returns immediately. Otherwise, it initializes $w_i^{(0)}:=\frac{nr_i}{R}$ and performs $N=\left\lceil p\log\left(\frac{8p\sqrt{R}}{\eps}\right)\right\rceil$ iterations, each of whose key step is the following, which we refer to as the ``Lewis update'':\[ w_i^{(t+1)}:=\blevsc{i}\left(\left(\bW^{(t)}\right)^{1/2-1/p}\bA\right), \text{ for all } i \in [k] \numberthis\label{eq:Lewis-update}.\]After all $N$ iterations, the algorithm returns $\what_i
:=
\left(
\frac{w_i^{(N)}}{(w_i^{(N-1)})^{1-2/p}}
\right)^{p/2}$,
for all $i\in[k].$ 

In words, \Cref{eq:Lewis-update} rescales the rows of $\bA$ according to the current block weights $\bw^{(t)}$, computes the exact row leverage scores of 
the rescaled matrix $\left(\bW^{(t)}\right)^{1/2-1/p}\bA$, sums them within each block, and replaces each block weight for the next iteration by the corresponding block leverage score. 
\Cref{eq:Lewis-update} is motivated by the defining fixed-point relation $\wbar_i=\blevsc{i}(\bWbar^{1/2-1/p}\bA)$ for every $i\in[k]$ for the true block Lewis weights $\bwbar$.
The same relation also motivates the final postprocessing:   since
$\wbar_i^{2/p}=\frac{\wbar_i}{\wbar_i^{1-2/p}}$,
the algorithm applies this same algebraic inversion to the final pair of iterates.


Our analysis centers on this fixed-point structure. Its key step is to show that one Lewis update contracts a KL divergence to $\bwbar$ by the factor $1-\frac{2}{p}$ (\Cref{sec:kl-decay}). Iterating this update brings $\bG(\bw)$ sufficiently close spectrally  to $\bG(\bwbar)$, and combining this spectral closeness with the postprocessing above yields  coordinatewise $(1\pm\eps)$-approximation to $\bwbar$ (\Cref{sec:coordinate-wise-control}).


\begin{algorithm}[t]
\caption{Fixed-point iteration}
\label{alg:main}
\begin{algorithmic}[1]
\Require A full-column-rank matrix $\bA\in\R^{m\times n}$  partitioned into $k$ row blocks $\bAblk{1}, \bAblk{2}, \dotsc, \bAblk{k}$ with $\bAblk{i}\in\R^{m_i\times n}$ for each $i\in[k]$, problem parameter $p>2$, and accuracy parameter $\eps\in(0, 1)$. 
\State Compute $r_i:=\rank(\bAblk{i})$ for all $i\in [k]$, and define $R:= \sum_{i = 1}^k r_i$. 
\If{$R=n$}
  \State \Return $\what_i=r_i$ for all $i\in[k]$. 
\EndIf
\State Initialize $w_i^{(0)}:= \frac{nr_i}{R}$ for all $i\in[k]$, and set $N = \left\lceil p \log\left(\frac{8p\sqrt{R}}{\eps}\right)\right\rceil$.
\For{$t=0,1,\ldots,N-1$}
  \State Set $\bW^{(t)}:=\Diag\left(w_1^{(t)}\bI_{m_1}, w_2^{(t)}\bI_{m_2}, \dotsc, w_k^{(t)}\bI_{m_k}\right)$.
  \State Compute $\levsc\left((\bW^{(t)})^{1/2-1/p}\bA\right)\in\R^m_{\geq0}$, the  exact row-wise leverage scores of  $(\bW^{(t)})^{1/2-1/p}\bA$. 
  \For{block $i=1,\ldots,k$}
    \State Set $w_i^{(t+1)}:= \blevsc{i}\left( (\bW^{(t)})^{1/2-1/p}\bA \right) $.
  \EndFor
\EndFor
\State \Return $\what_i:=\left(\frac{w_i^{(N)}}{(w_i^{(N-1)})^{1-2/p}}\right)^{p/2}$ for all $i\in [k]$.
\end{algorithmic}
\end{algorithm}

\paragraph{Notation.}
For the analysis in this section and in \Cref{sec:intuition}, we use the following notation. We set $\omrp:=1-\frac{2}{p}\in(0,1)$ and define  $\Omega:=\{\bu\in\Rkpos:\one^\top\bu=n\}$, the set block weights live in.

More generally, for positive vectors $\bx,\by\in\R_{>0}^d$ of the same dimension satisfying $\one^\top\bx=\one^\top\by=n$, we define the scaled Kullback-Leibler (KL) divergence
\[
\kldiv{\bx}{\by}
:=
\sum_{i=1}^d x_i\log\left(\frac{x_i}{y_i}\right).
\numberthis\label{eq:notation-KLdiv}
\]
For probability distributions $\mu$ and $\nu$, we use $\KL{\mu}{\nu}$ for the usual KL divergence.

Given $\bw\in\Omega$ and a specific block partitioning, let $\bW$ denote its block weight matrix as defined in \Cref{sec:intro-notationPrelims}, and define  $\bG(\bw):=\bA^\top\bW^\omrp\bA$. Since $\bA$ has full column rank, $\bG(\bw)\succ\bzero$.

We use $\calT:\R_{>0}^k\to\Rkpos$ as shorthand for the main update in \Cref{alg:main}, namely
\[
\calT_i(\bw)
:=
\blevsc{i}\left(\bW^{\omrp/2}\bA\right)
=
w_i^\omrp
\Tr\left(
\bAblk{i}
\left(\bA^\top\bW^\omrp\bA\right)^{-1}
\bAblk{i}^\top
\right) = w_i^\omrp \Tr\left(\bAblk{i}\bG(\bw)^{-1}\bAblk{i}^\top\right)
\text{ for all }i\in[k].
\]
Thus, \Cref{eq:Lewis-update} can be written as $\bw^{(t+1)}=\calT(\bw^{(t)})$, a formulation we often use. By \Cref{fact:block-levsc} and the assumption that every block is nonzero, $\calT(\Omega)\subseteq\Omega$. Together with $\bw^{(0)}\in\Omega$, this implies $\bw^{(t)}\in\Omega$ for every iterate of \Cref{alg:main}. The true block Lewis weights likewise satisfy $\bwbar=\calT(\bwbar)\in\Omega$ by \Cref{def:blockLW} and \Cref{fact:block-levsc}.
When we specialize to row-wise Lewis weights in \Cref{sec:intuition}, we overload $\Omega$ with $\{\bu\in\Rmspos:\one^\top\bu=n\}$ and $\calT$ for the row-wise Lewis map $\calT_i(\bw):=\levsc{i}\left(\bW^{\omrp/2}\bA\right)$.

\subsection{The block Lewis update contracts KL divergence}
\label{sec:kl-decay}

\begin{lemma}
\label{lem:kl-contraction}
Let  $\bw\in\Omega$. Then  $\kldiv{\bwbar}{\calT(\bw)}
    \leq
    \omrp\kldiv{\bwbar}{\bw}.$
\end{lemma}

\begin{proof}
By  the definition of block leverage scores and $\bwbar=\calT(\bwbar)$, we have \[ \calT_i(\bw)=w_i^\omrp \Tr\left(\bAblk{i}\bG(\bw)^{-1}\bAblk{i}^\top\right), \qquad \wbar_i = \wbar_i^{\omrp} \Tr\left(\bAblk{i}\bG(\bwbar)^{-1}\bAblk{i}^\top\right).\] Therefore, by expanding out the definition in \Cref{eq:notation-KLdiv}, we have \[\kldiv{\bwbar}{\calT(\bw)} = \omrp \kldiv{\bwbar}{\bw} + \sum_{i = 1}^k \wbar_i \log\left(\frac{\Tr\left(\bAblk{i}\bG(\bwbar)^{-1}\bAblk{i}^\top\right)}{\Tr\left(\bAblk{i}\bG(\bw)^{-1}\bAblk{i}^\top\right)}\right). \numberthis\label{eq:KLcon-2}\] We  complete the  proof  by showing next that the second term of \Cref{eq:KLcon-2} is nonpositive. To this end, we first note that by Jensen's inequality applied to $x\mapsto\log(x)$ and the fact that $\sum_{i=1}^k \wbar_i = n$, we have  \[ \sum_{i = 1}^k \frac{\wbar_i}{n} \log\left(\frac{\Tr\left(\bAblk{i}\bG(\bwbar)^{-1}\bAblk{i}^\top\right)}{\Tr\left(\bAblk{i}\bG(\bw)^{-1}\bAblk{i}^\top\right)}\right) \leq \log\left( \sum_{i=1}^k \frac{\wbar_i}{n} \frac{\Tr\left(\bAblk{i}\bG(\bwbar)^{-1}\bAblk{i}^\top\right)}{\Tr\left(\bAblk{i}\bG(\bw)^{-1}\bAblk{i}^\top\right)} \right),\] and it  suffices to show that \[\sum_{i=1}^k {\wbar_i} \frac{\Tr\left(\bAblk{i}\bG(\bwbar)^{-1}\bAblk{i}^\top\right)}{\Tr\left(\bAblk{i}\bG(\bw)^{-1}\bAblk{i}^\top\right)} \stackrel{?}{\leq} n.\numberthis\label{eq:KLcon-3}\]To this end, we use $\bwbar=\calT(\bwbar)$ and apply matrix Cauchy-Schwarz inequality $\left|\Tr\left(\bX^\top\bY\right)\right|\leq \|\bX\|_{\mathrm{F}}\|\bY\|_{\mathrm{F}}$ with $\bX:=\bG(\bw)^{1/2}\bG(\bwbar)^{-1}\bAblk{i}^\top$ and $\bY:= \bG(\bw)^{-1/2}\bAblk{i}^\top$ and first obtain \[\wbar_i \frac{\Tr\left(\bAblk{i} \bG(\bwbar)^{-1}\bAblk{i}^\top\right)}{\Tr\left(\bAblk{i} \bG(\bw)^{-1}\bAblk{i}^\top\right)} = \wbar_i^{\omrp} \frac{\left[\Tr \left(\bAblk{i}\bG(\bwbar)^{-1}\bAblk{i}^\top\right)\right]^2}{\Tr\left(\bAblk{i}\bG(\bw)^{-1}\bAblk{i}^\top\right)} \leq \wbar_i^{\omrp} \Tr\left(\bAblk{i} \bG(\bwbar)^{-1}\bG(\bw)  \bG(\bwbar)^{-1} \bAblk{i}^\top\right).\] Summing the above inequality over $i = 1, 2, \dotsc, k$ and applying the definition of $\bG(\bwbar)$ and $\bG(\bw)$, linearity of trace, and the invariance of trace to cyclic permutation, we have
\begin{align*}
\sum_{i=1}^k \wbar_i\frac{\Tr\left(\bAblk{i}\bG(\bwbar)^{-1}\bAblk{i}^\top\right)}{\Tr\left(\bAblk{i}\bG(\bw)^{-1}\bAblk{i}^\top\right)} &\leq \sum_{i=1}^k \wbar_i^\omrp \Tr\left(\bAblk{i}\bG(\bwbar)^{-1}\bG(\bw) \bG(\bwbar)^{-1} \bAblk{i}^\top\right)= \Tr\left(\bG(\bw)\bG(\bwbar)^{-1}\right).\numberthis\label{eq:KLcon-4}
\end{align*}
Next, again by $\bwbar=\calT(\bwbar)$, H\"older's inequality with exponents $\frac{1}{\omrp}$ and $\frac{1}{1-\omrp}$ (valid since we restrict ourselves to $p>2$), and the fact that $\bw\in\Omega$, we have \[ \Tr\left(\bG(\bw)\bG(\bwbar)^{-1}\right)
= \sum_{i=1}^k w_i^\omrp \Tr\left( \bAblk{i}\bG(\bwbar)^{-1} \bAblk{i}^\top\right) = \sum_{i=1}^k w_i^\omrp \wbar_i^{1-\omrp}\leq \left(\sum_{i=1}^k w_i\right)^\omrp \left(\sum_{i=1}^k \wbar_i\right)^{1-\omrp} = n. \numberthis\label{eq:KLcon-5}\]  Combining \Cref{eq:KLcon-4} and \Cref{eq:KLcon-5} completes the proof of \Cref{eq:KLcon-3}, which finishes the proof of the claim. 
\end{proof}

\begin{corollary}\label{cor:kl-decay}Assume $R>n$. 
Then for every $t\geq0$, the iterates  of \Cref{alg:main} satisfy \[\kldiv{\bwbar}{\bw^{(t)}}\leq\omrp^t n \log\left(\frac{R}{n}\right).\]
\end{corollary}
\begin{proof}
The assumption $R>n$ implies \Cref{alg:main} does not terminate at Line 3 and instead initializes $w_i^{(0)} = \frac{nr_i}{R}$. We also know from \Cref{fact:block-levsc} that $0\leq \wbar_i \leq r_i$ and $\sum_{i=1}^k \wbar_i = n$. Therefore, \[\kldiv{\bwbar}{\bw^{(0)}}= \sum_{i=1}^k \wbar_i\log\left(\frac{\wbar_i}{w^{(0)}_i}\right)\leq n\log\left(\frac{R}{n}\right). \numberthis\label{eq:cor-con-1}\]Applying \Cref{lem:kl-contraction} at $\bw = \bw^{(s)}$  gives $\kldiv{\bwbar}{\bw^{(s+1)}}= \kldiv{\bwbar}{\calT(\bw^{(s)})}\leq \omrp\kldiv{\bwbar}{\bw^{(s)}}$; repeatedly applying this at $s = 0, 1, 2, \dotsc, t-1$ and plugging in \Cref{eq:cor-con-1}  gives the claimed decay. 
\end{proof} 

\subsection{From KL contraction to a spectral bound}
\label{sec:coordinate-wise-control}
 For two compatible matrices $\bX, \bY\succ\bzero$, we define the Bregman divergence~\cite{bregman1967relaxation} \[\logdetdiv{\bX}{\bY}:= -\log\det\bX+\log\det\bY + \Tr\left(\bY^{-1}\left(\bX-\bY\right)\right).\] 
\begin{lemma}\label{lem:normalized-gram-eigs-small}
Let $\bw\in \Omega$. Then, 
 we have $\logdetdiv{\bG(\bw)}{\bG(\bwbar)} \leq \omrp\kldiv{\bwbar}{\bw}$. 
\end{lemma}
\begin{proof}
Our proof strategy is as follows: We will show that \[\Tr\left(\bG(\bw)\bG(\bwbar)^{-1}\right)\stackrel{?}{\leq} n, \qquad \log\det(\bG(\bwbar)) - \log\det(\bG(\bw)) \stackrel{?}{\leq}\omrp\kldiv{\bwbar}{\bw}.\numberthis\label{eq:spectral-1}\] Using these bounds in the definition of $\logdetdiv{\bG(\bw)}{\bG(\bwbar)}$ gives the below chain of steps: \begin{align*}\logdetdiv{\bG(\bw)}{\bG(\bwbar)} &= \Tr\left(\bG(\bw)\bG(\bwbar)^{-1}\right) - n + \log\det(\bG(\bwbar)) - \log\det(\bG(\bw))\\ & \leq \omrp\kldiv{\bwbar}{\bw}.\end{align*} which  concludes the proof. We now show \Cref{eq:spectral-1}. We already have  $\Tr\left(\bG(\bw)\bG(\bwbar)^{-1}\right)\leq n$  from \Cref{eq:KLcon-5}. To prove the second inequality in \Cref{eq:spectral-1}, we start with the fact that $\bx\mapsto\log\det\bG(e^{\bx})$ is convex (cf. \Cref{lem:logdetGexp-is-cvx}). Then, by the linear underestimator property of convex functions and the identities $\frac{\partial}{\partial x_i}\log\det\bG\left(e^{\bx}\right) = \omrp \calT_{i}\left(e^{\bx}\right)$ (cf. \Cref{lem:logdetGexp-is-cvx}) and $\bwbar=\calT(\bwbar)$,  we have \begin{align*}
\log\det\bG(\bw) &= \log\det\bG\left(e^{\log \bw}\right)\\ &\geq \log\det\bG\left(e^{\log \bwbar}\right) + \nabla_{\bx}\log\det\bG\left(e^{\bx}\right)|_{\bx=\log\bwbar}^\top (\log \bw-\log\bwbar)\\  
&= \log\det\bG(\bwbar) + \omrp\cdot \sum_{i = 1}^k \wbar_i \log\left(\frac{w_i}{\wbar_i}\right) = \log\det\bG(\bwbar) - \omrp \kldiv{\bwbar}{\bw}. 
\end{align*} Rearranging the terms  proves the second inequality in \Cref{eq:spectral-1}, thereby proving the lemma. 
\end{proof}

\begin{lemma}
\label{lem:coordinatewise-approx}
Let $\bw\in\Omega$ satisfy $\logdetdiv{\bG(\bw)}{\bG(\bwbar)}\leq \frac{\eps^2}{64p^2}$. Define $\what_i := \left[\Tr\left(\bAblk{i}\bG(\bw)^{-1}\bAblk{i}^\top\right)\right]^{\frac{p}{2}}$ for all $i \in [k]$. Then, $(1-\eps)\wbar_i \leq \what_i \leq (1+\eps)\wbar_i$ for all $i \in [k]$.
\end{lemma}
\begin{proof}
Set $\delta := \frac{\eps^2}{64p^2}$ and $a:= 2\sqrt{\delta}$. Then, under the premise of our lemma (and using our standing assumption that $p>2$), \cite[Lemma~20]{pmlr-v336-gribling26a} yields: \[(1-a)\bG(\bwbar)\preceq \bG(\bw)\preceq (1+a) \bG(\bwbar).\] Taking inverses, pre- and post-multiplying  by appropriate matrices, and taking the trace yields: \[\frac{1}{1+a}\Tr\left(\bAblk{i}\bG(\bwbar)^{-1}\bAblk{i}^\top\right)\leq \Tr\left(\bAblk{i}\bG(\bw)^{-1}\bAblk{i}^\top\right)\leq \frac{1}{1-a}\Tr\left(\bAblk{i}\bG(\bw)^{-1}\bAblk{i}^\top\right).\numberthis\label{eq:coord1}\] Since $\bwbar= \calT(\bwbar)$, we have $\Tr\left(\bAblk{i}\bG(\bwbar)^{-1}\bAblk{i}^\top\right) = \wbar_i^{2/p}.$ Substituting this back into \Cref{eq:coord1}, exponentiating by $p/2$, and plugging in the definitions of $\bwbar$ and $\widehat{\bw}$ gives
\[(1+a)^{-p/2}\wbar_i\leq \what_i \leq (1-a)^{-p/2}\wbar_i.\numberthis\label{eq:coord2}\]
We simplify both sides of the above inequality chain, one at a time. First, we have $(1+a)^{-p/2} \geq e^{-ap/2} = e^{-\varepsilon/8}\geq 1-\varepsilon/8\geq 1-\varepsilon.$ Since $a<1/2$, we have $-\log(1-a)\leq 2a$, and therefore $(1-a)^{-p/2} \leq e^{ap} =e^{\eps/4} \leq 1+\eps$.  These two simplifications, plugged into \Cref{eq:coord2} yield the claim.  
\end{proof}

\subsection{Proof of \Cref{thm:blkLW}}\label{sec:main-proof}
\begin{proof}[Proof of \Cref{thm:blkLW}]
First, we note that by subadditivity of rank under vertical concatenation, it is always  the case that $R:= \sum_{i=1}^k \rank\left(\bAblk{i}\right)\geq \rank(\bA)=: n$. 
When $R = n$, one can infer from \Cref{fact:block-levsc} that each $\blevsc{i}=\rank\left(\bAblk{i}\right)$ for all $i\in[k]$, and, as designed in \Cref{alg:main}, the block Lewis weights are  $\what_i=\rank\left(\bAblk{i}\right)$ for all $i\in[k]$, thereby proving correctness in this case. The claimed iteration bound is immediate in this case, since \Cref{alg:main} computes no leverage scores.  
We now assume the other case, $R>n$.
By \Cref{lem:normalized-gram-eigs-small} applied to $\bw= \bw^{(N-1)}$, \Cref{cor:kl-decay} applied to $t=N-1$, the choice of $\omrp=1-\frac{2}{p}$, and the choice of $N = \left\lceil p\log\left(\frac{8p\sqrt{R}}{\eps}\right)\right\rceil$ in \Cref{alg:main}, we have 
\[ \logdetdiv{\bG(\bw^{(N-1)})}{\bG(\bwbar)}\leq \omrp \kldiv{\bwbar}{\bw^{(N-1)}}\leq \gamma^N n \log\left(\frac{R}{n}\right)\leq R e^{-2N/p}\leq \frac{\eps^2}{64p^2}.\]\looseness=-1 Thus, the premise of \Cref{lem:coordinatewise-approx} is satisfied at $\bw=\bw^{(N-1)}$; further, the output of \Cref{alg:main} is exactly $\bwhat$ from \Cref{lem:coordinatewise-approx} with $\bw=\bw^{(N-1)}$; hence, \Cref{lem:coordinatewise-approx} implies $(1-\eps)\wbar_i \leq \what_i\leq (1+\eps)\wbar_i$ for all $i\in [k]$ for the output of \Cref{alg:main}. Each of the $N$ iterations of \Cref{alg:main} computes  one exact  leverage score vector and  sums the entries within each block, implying the claimed  complexity. 
\end{proof}

\section{Two additional perspectives on the Lewis fixed-point iteration}\label{sec:intuition}
\looseness=-1By \Cref{sec:main-results}, the update in \Cref{eq:Lewis-update} contracts the KL divergence to the true Lewis weights by $1-\frac{2}{p}$ at each iteration, driving the $\widetilde{O}(p)$ rate in \Cref{thm:blkLW}. We give two additional perspectives on this iteration.

\subsection{Global dynamics: a volume sampling perspective}\label{sec:global-intuition}

First, we present the volume sampling perspective on the \emph{row-wise} Lewis update $\bw\mapsto\calT(\bw)$. The main idea is: 
Under the external tilt $\bw^{1-2/p}$, the marginal inclusion probability vector of the resulting volume sampling distribution is exactly $\calT(\bw)$. 
Combining this with entropic independence~\cite{anari2022entropic} gives an alternate proof of the KL contraction in \Cref{lem:kl-contraction} and explains the appearance of the factor $1-\frac{2}{p}$.
Below, we elaborate on this argument in the row-wise setting; the same argument extends to blocks by grouping the row marginals within each block.

  For any $\bv\in\Rmspos$, let $\pi_{\bv}$ denote the externally tilted volume sampling distribution 
  over $n$-element subsets of the $m$ rows of $\bA$:
  \[\pi_{\bv}(S) := \frac{(\det(\bA_S))^2\prod_{i \in S} v_i}{Z(\bv)},\, S\subseteq[m], |S| = n; \quad Z(\bv) := \sum_{\substack{S\subseteq[m],\\|S|=n}} \det(\bA_S)^2\prod_{i\in S}v_i.\numberthis\label{eq:EI-con-0}\] 
Since $\det(\bA_S)^2\prod_{i\in S}v_i = \det\left(\left(\Diag(\sqrt{\bv})\bA\right)_S\right)^2$, one can equivalently describe $\pi_{\bv}$ as the volume sampling distribution induced by the matrix $\Diag(\sqrt{\bv})\bA$. For this $n$-element volume sampling distribution, it is known that the marginal probability that row $i$ is included equals the $i^\mathrm{th}$ leverage score of $\Diag(\sqrt{\bv})\bA$~\cite[Proposition 4]{derezinski2018reverse}\cite[Theorem 8]{derezinski2021determinantal}, i.e., \[ \Prob_{S\sim \pi_{\bv}}\left(i\in S\right) = \levsc{i}\left(\Diag\left(\sqrt{\bv}\right)\bA\right), \, \text{ for all } i\in[m]. \numberthis\label{eq:EI-con-01}\] We now specialize \Cref{eq:EI-con-01} to $\bv=\bw^{\omrp}$ and $\bv=\bwbar^\omrp$ and use $\Diag\left(\sqrt{\bw^\omrp}\right)\bA = \bW^{\omrp/2}\bA$ and the fixed point characterization $\bwbar=\calT(\bwbar)$  to conclude the following facts we shortly use: \begin{equation}\begin{aligned}\Prob_{S \sim \pi_{\bw^\omrp}}\left(i\in S\right) &= \levsc{i}\left(\bW^{\omrp/2}\bA\right) = \calT_i(\bw), \, \text{ for all } i\in[m]. \\
\Prob_{S\sim\pi_{\bwbar^\omrp}}\left(i\in S\right) &= \bwbar_i, \, \text{ for all } i\in[m]. \label{eq:EI-con-02}\end{aligned}\end{equation} 
We are now ready to use the above established facts. 
To this end, we invoke entropic indepdence~\cite{anari2022entropic}, a strong
data-processing property for distributions on fixed-size sets: Let $\mu$ and $\nu$ be two distributions on $n$-element subsets of $[m]$, with marginal inclusion probability vectors $\ba$ and $\bb$, respectively. If, after sampling a set, we retain only one uniformly chosen random element, the resulting distributions on $[m]$ 
are $\frac{\ba}{n}$ and $\frac{\bb}{n}$. The standard data-processing inequality  implies that this operation cannot increase KL divergence; it was shown~\cite{anari2022entropic} that if, however, $\nu$ is entropically independent with parameter $\alpha=1$, the following stronger bound holds: 
\[ 
\KL{\frac{\ba}{n}}{\frac{\bb}{n}} \leq \frac{1}{n}\KL{\mu}{\nu}. 
\]
 Equivalently, using
\Cref{eq:notation-KLdiv} on the left-hand side, we have
\[
\kldiv{\ba}{\bb}\leq \KL{\mu}{\nu}.
\numberthis\label{eq:EI-con-1}
\]In our context, fix $\bw\in\Omega$ and set
$\mu=\pi_{\bwbar^\omrp}$ and $\nu=\pi_{\bw^\omrp}$.
The distribution $\nu=\pi_{\bw^\omrp}$ is a determinantal volume sampling measure and hence  entropically independent with parameter $\alpha=1$~\cite{anari2022entropic,anari2025optimal}.
Therefore, \Cref{eq:EI-con-1} applies.
By \Cref{eq:EI-con-02}, the marginal vectors of $\mu$ and $\nu$ are $\bwbar$ and $\calT(\bw)$, respectively, and hence
\[
\kldiv{\bwbar}{\calT(\bw)}
\leq
\KL{\pi_{\bwbar^\omrp}}{\pi_{\bw^\omrp}}.
\numberthis\label{eq:EI-con-2}
\] We now expand the right-hand side above: \begin{align*}\KL{\pi_{\bwbar^\omrp}}{\pi_{\bw^\omrp}} &=  \E_{S\sim \pi_{\bwbar^\omrp}}\left(\log\left(\frac{\pi_{\bwbar^\omrp}(S)}{\pi_{\bw^\omrp}(S)}\right)\right)  \\  
&= \omrp\E_{S\sim \pi_{\bwbar^\omrp}} \left( \sum_{i \in S}\log\left(\frac{\wbar_i}{w_i}\right)\right)  + \log\left(\frac{Z(\bw^\omrp)}{Z(\bwbar^\omrp)}\right) \\ 
&= \omrp\sum_{i=1}^m \Prob_{S\sim \pi_{\bwbar^\omrp}}(i\in S) \log\left(\frac{\wbar_i}{w_i}\right) + \log\left(\frac{Z(\bw^\omrp)}{Z(\bwbar^\omrp)}\right) \\ 
&= \gamma\sum_{i=1}^m \wbar_i \log\left(\frac{\wbar_i}{w_i}\right) + \log\left(\frac{Z(\bw^\omrp)}{Z(\bwbar^\omrp)}\right) \\
&= \gamma \kldiv{\bwbar}{\bw} +\log\det\bG(\bw)-\log\det\bG(\bwbar),\numberthis\label{eq:EI-con-3}
\end{align*} where the first step is by the definition of the usual KL divergence, the second step by \Cref{eq:EI-con-0}, the third step by representing membership in $S$ via  indicator variables and using linearity of expectations, the fourth step by \Cref{eq:EI-con-02}, and the fifth step uses Cauchy-Binet \Cref{{eq:logdetG-1}} and our definition of $\bG$ to conclude $Z(\bw^\omrp) = \det\left(\bA^\top\bW^\omrp\bA\right)=\det(\bG(\bw))$. Finally, by the variational characterization of Lewis weights~\cite{cohen2015lp,lee2019solving}, we have $\bwbar\in\arg\max_{\bu\in\Omega}\log\det\bG(\bu),$ which simplifies 
\Cref{eq:EI-con-3} to \[ \KL{\pi_{\bwbar^\omrp}}{\pi_{\bw^\omrp}} \leq \omrp \kldiv{\bwbar}{\bw}.\numberthis\label{eq:EI-con-4}\] Chaining together \Cref{eq:EI-con-2} and \Cref{eq:EI-con-4}  yields $\kldiv{\bwbar}{\calT(\bw)}\leq \left(1-\frac{2}{p}\right)\kldiv{\bwbar}{\bw}$, thereby recovering the  contraction in \Cref{lem:kl-contraction} in the row-wise setting. 

Thus, the above calculations explain the $1-\frac{2}{p}$ factor: the \Cref{eq:Lewis-update} corresponds to the external-field tilt $\bw^{\omrp}$, and taking the log-likelihood ratio between the two tilted volume sampling distributions pulls this exponent directly outside their KL divergence. 

\subsection{Local dynamics: linearization at the fixed point}
\label{sec:local-dyn}
The above volume sampling perspective explains the \emph{global} KL contraction in \Cref{lem:kl-contraction} and the associated $1-\frac{2}{p}$ factor; however, this contraction does not characterize the asymptotic convergence rate of the iterates near the fixed point $\bwbar$. Motivated by stable geometric decay we observe in exploratory numerical experiments, we develop a more precise theory of these \emph{local} dynamics. 
To this end, we linearize the row-wise \Cref{eq:Lewis-update} around $\bwbar$; the spectral radius of the resulting Jacobian predicts the asymptotic local contraction rate. 
As seen in \Cref{fig:local-dyn}, this predicted rate is closely matched by the observed dynamics on independently generated test instances.

\begin{lemma}\label{lem:local-dyn} Let  $\wbar\in\Rmspos$ be the row-wise Lewis weights of $\bA$, and let $\calT(\bw):=\levsc\left(\bW^{\omrp/2}\bA\right)$ denote \Cref{eq:Lewis-update} applied row-wise. Define the  matrices $\bP:= \bWbar^{\omrp/2}\bA\left(\bA^\top \bWbar^\omrp \bA\right)^{-1}\bA^\top\bWbar^{\omrp/2}\in\R^{m\times m}$, and $\bK:= \bWbar^{-1}(\bP\circ\bP) \in\R^{m \times m}$. Define the Jacobian $\bJ:= \left.\left[\frac{\partial}{\partial x_j}\log\calT_i(e^{\bx})\right]\right|_{i, j\in [m]}\lvert_{\bx = \log\bwbar}$. Then, $\bJ = \omrp(\bI-\bK).$ Moreover, all eigenvalues of $\bK$ lie in $[0,1]$. Therefore, the spectral radius of $\bJ$ is
\[
\mu_{\mathrm{pred}}
=
\left(1-\frac{2}{p}\right)
\left(1-\lambda_{\min}(\bK)\right).
\]
\end{lemma}
\begin{proof}
For this proof, we introduce the shorthand notation  $\bq(\bx):= e^{\omrp \bx/2}$,   $\bQ(\bx):= \Diag(\bq(\bx))$, and $\bP(\bx):= \bQ(\bx)\bA\left(\bA^\top\bQ(\bx)^2\bA\right)^{-1}\bA^\top \bQ(\bx)$. Then, it follows that \[\diag(\bP(\bx))= \levsc(\bQ(\bx)\bA) = \calT(e^{\bx}).\numberthis\label{eq:locdyn-0}\] Fix $\bx\in\R^m$ so that $\bq(\bx)\in\R^m$ is fixed too. Then, specializing \cite[Lemma~49]{lee2019solving} to $\bQ(\bx)\bA\in\R^{m\times n}$ and using \Cref{eq:locdyn-0} gives, for any perturbation $\bu\in\R^m$, \[\frac{d}{dt} \levsc\left(\Diag\left(\bq(\bx) + t\bu\right)\bA\right)|_{t=0} = 2\left(\Diag(\calT(e^{\bx})) - \bP(\bx)\circ\bP(\bx) \right)\bQ(\bx)^{-1}\bu.\]Next, perturbing $\bx$ in the direction $\bh$ gives $\frac{d}{dt}\bq(\bx+t\bh)|_{t=0} = \frac{\omrp}{2}\bQ(\bx)\bh.$ 
Therefore,  chain rule gives \begin{align*}\frac{d}{dt}\calT(e^{\bx+t\bh})|_{t=0} &= 2\left(\Diag(\calT(e^{\bx})) - \bP(\bx)\circ\bP(\bx)\right)\bQ(\bx)^{-1}\left(\frac{\omrp}{2}\bQ(\bx)\bh\right)\\ 
&= \omrp\cdot\left(\Diag(\calT(e^{\bx})) - \bP(\bx)\circ\bP(\bx)\right)\bh. 
\end{align*} Again, by chain rule, we have \begin{align*}
\frac{d}{dt}\log\calT(e^{\bx+t\bh})|_{t=0} &= \Diag(\calT(e^{\bx}))^{-1}\frac{d}{dt}\calT(e^{\bx+t\bh})|_{t=0}\\ 
&= \omrp\left(\bI- \Diag(\calT(e^{\bx}))^{-1}\left(\bP(\bx)\circ\bP(\bx)\right)\right)\bh.\numberthis\label{eq:locdyn-2}\end{align*} Setting $\bx=\log\bwbar$ and using the fact that $\bwbar=\calT(\bwbar)$ implies $\bP(\log\bwbar)= \bP$; plugging these, along with the definition of $\bK$, in \Cref{eq:locdyn-2} gives the claimed Jacobian. 

Next, we bound the spectrum of $\bK$. By Part (2) of \cite[Lemma~47]{lee2019solving}, we have 
$
    \bzero
    \preceq
    \bP\circ\bP
    \preceq
    \bWbar.
$
Pre and post multiplying by $\bWbar^{-1/2}$, this simplifies to 
\[
    \bzero
    \preceq
    \bWbar^{-1/2}
    (\bP\circ\bP)
    \bWbar^{-1/2}
    \preceq
    \bI.
\] Finally, note that $\bK$ is similar to the above $\bWbar^{-1/2}(\bP\circ\bP)\bWbar^{-1/2}$, thus implying $0\leq \lambda_i(\bK)\leq 1$ for all $i = 1, 2, \dotsc, m$. This immediately yields the claimed spectral radius, thus finishing the proof. 
\end{proof}

\noindent Thus, unlike the global KL factor $1-\frac{2}{p}$, the local contraction rate depends on the local geometry and can be strictly faster. We now explain how we test this theory in our experiments. 

\paragraph{Connection to the measured residual.}
To connect $\mu_{\mathrm{pred}}$ from \Cref{lem:local-dyn} to the quantity measured in \Cref{sec:expts-A0}, let $\bx^{(t)}:=\log\bw^{(t)}$ and $\bxbar:=\log\bwbar$. 
Rather than measuring the error to the fixed point directly, we use the residual 
\[
R_t
:=
\max_{i\in[m]}
\left|
\log\frac{w_i^{(t+1)}}{w_i^{(t)}}
\right|
=
\left\|
\bx^{(t+1)}-\bx^{(t)}
\right\|_\infty.
\]
By the definition of $\bJ$ in \Cref{lem:local-dyn} and a Taylor expansion of the map $\bx\mapsto\log\calT(e^{\bx})$ around $\bxbar$, as the iterates approach the fixed point we have
\[
\bx^{(t+2)}-\bx^{(t+1)}
=
\bJ\left(\bx^{(t+1)}-\bx^{(t)}\right)
+
o\left(
\left\|\bx^{(t+1)}-\bx^{(t)}\right\|
\right).
\]Thus, the Jacobian $\bJ$ from \Cref{lem:local-dyn} controls the local decay of our residual.
In particular, when the local trajectory has a nonzero component in an eigenspace associated with $\rho(\bJ)$, we expect $
\frac{R_{t+1}}{R_t}
\longrightarrow
\rho(\bJ)
=
\mu_{\mathrm{pred}}.$ 
Equivalently, in the local regime,
$
R_t \approx C\mu_{\mathrm{pred}}^t,
$ 
so that $\log R_t$ is approximately linear in $t$ with slope $\log\mu_{\mathrm{pred}}$. 
In \Cref{sec:expts-A0}, we estimate the observed factor $\mu_{\mathrm{obs}}$ from this decay and compare it directly with $\mu_{\mathrm{pred}}$ (see \Cref{fig:local-dyn}). Thus, \Cref{lem:local-dyn} yields a directly observable prediction for the asymptotic decay of our residual.

\section{Numerical experiments}
\label{sec:expts}
We implemented our algorithm and tested it on synthetic instances and real data. 

\subsection{Synthetic experiments}\label{sec:syn_expt}
We begin with a question about local dynamics: we ask if the asymptotic contraction observed near the Lewis weight fixed point agrees with the theory developed in \Cref{sec:local-dyn}.
We then ask whether the global $\widetilde O(p)$ dependence in \Cref{cor:ordinary-main} is reflected in the observed iteration counts of the ordinary Lewis weight iteration. 
We first study this question on i.i.d.\ Gaussian matrices, and then check whether the same qualitative behaviour persists under two departures from this baseline: first on a structured, well-conditioned matrix family with rows at two significantly different scales, and then on ill-conditioned matrices designed to stress the numerical implementation.
Finally, we turn from row-wise to block Lewis weights and ask whether the block Lewis weight iteration exhibits the analogous qualitative dependence on $p$ suggested by \Cref{thm:blkLW}.

\subsubsection{Local dynamics}
\label{sec:expts-A0}

To test the local prediction in \Cref{sec:local-dyn}, we generate $30$ independent matrices $\bA\in\R^{100\times20}$ with i.i.d.\ standard Gaussian entries and repeatedly apply the row-wise specialization of \Cref{eq:Lewis-update} for $14$ values of $p$ between $3$ and $64$. We measure progress by
$R_t:=\max_{i\in[m]}\left|\log\left(\frac{w_i^{(t+1)}}{w_i^{(t)}}\right)\right|$,
the maximum coordinatewise change between successive iterates; thus $R_t=0$ exactly at a fixed point, while $R_t>0$ quantifies the maximum coordinatewise violation of the fixed-point relation.

For each matrix and $p$, \Cref{lem:local-dyn} predicts the asymptotic contraction factor
$\mu_{\mathrm{pred}}=\left(1-\frac{2}{p}\right)(1-\lambda_{\min}(\bK))$.
We estimate the observed factor $\mu_{\mathrm{obs}}$ by fitting $\log R_t$ against $t$ over the prespecified window $10^{-8}\leq R_t\leq10^{-4}$, chosen to probe this local regime while avoiding the floating-point noise floor. This gives $30\times14=420$ prediction--observation comparisons.

\begin{figure}[!htbp]
    \centering
    \includegraphics[width=0.95\textwidth]{paperFigures/A0_main_agreement_figure.pdf}
    \caption{Predicted and observed local contraction factors for our Lewis weight iteration on Gaussian matrices.}
    \label{fig:local-dyn}
\end{figure}

The predicted and observed contraction factors closely agree across the tested range of $p$ (\Cref{fig:local-dyn}). In $418$ of the $420$ runs, a straight-line fit to $\log R_t$ over the primary window explains at least $99\%$ of its variation, consistent with geometric decay. The remaining two runs, at $p=61$ and $p=64$ for the same matrix, reach a floating-point noise floor near $10^{-8}$; over the prespecified earlier window $10^{-7}\leq R_t\leq10^{-4}$, all $420$ runs satisfy the same $99\%$ criterion. Thus the observed local dynamics provide empirical support for the prediction in \Cref{lem:local-dyn}.

\subsubsection{{Row-wise Lewis weights: iteration count and robustness}}\label{sec:A1-A2-A4}
\looseness=-1We now study the iteration counts and numerical robustness of the row-wise Lewis update. 


\paragraph{Iteration count of row-wise Lewis weights.}
We first ask how the row-wise iteration count varies with $p$ on generic dense instances. We generate 50 fixed matrices $\bA\in\R^{100\times20}$ with i.i.d.\ standard Gaussian entries and reuse them across an irregular $p$-grid from $3$ to $127.3$. Starting from the uniform initialization, we use $R_t:=\max_i\left|\log\frac{w_i^{(t+1)}}{w_i^{(t)}}\right|<10^{-8}$ as our primary stopping rule, and also record the first crossing of $10^{-4}$ and $10^{-6}$. For each $p$, the left panel of \Cref{fig:A1} summarizes the $10^{-8}$ stopping time across the 50 matrices by its median and empirical 20th and 80th percentiles, while the right reports the median stopping time at all three thresholds.

\begin{figure}[!htbp]
    \centering
    \includegraphics[width=0.95\textwidth]{paperFigures/A1_combined_panels.pdf}
    \caption{Row-wise Lewis weight round scaling with $p$ (left) and sensitivity to the stopping threshold (right). }
    \label{fig:A1}
\end{figure}


Across the tested $p$, the iteration count grows approximately linearly with $p$, varying little across instances; this qualitative behaviour persists at thresholds $10^{-4}$ and $10^{-6}$ (\Cref{fig:A1}). 
Thus, at fixed residual thresholds, our experiments are consistent with the $\widetilde O(p)$ dependence in \Cref{cor:ordinary-main}.

\paragraph{Structured row geometry.}
We next ask whether the behaviour above is peculiar to generic Gaussian matrices.
For each of 50 instances, we construct
$\bA =
    \begin{bmatrix}
        \bI_{20}\\[1mm]
        \frac{1}{4} \widetilde{\bG}
    \end{bmatrix}
    \in \R^{100\times 20}$,
where the 80 rows of $\widetilde{\bG}$ are independent Gaussian directions normalized to unit Euclidean norm. Thus the 20 ``anchor'' rows have norm one and the 80 ``cloud'' rows norm $\frac14$. This creates leverage heterogeneity without poor conditioning: the median anchor/cloud $p=2$ leverage ratio is about $16.8$, while the spectral condition number
$\kappa_2(\bA):=\frac{\sigma_{\max}(\bA)}{\sigma_{\min}(\bA)}$
ranges from roughly $1.16$ to $1.21$; see the right panel of \Cref{fig:structured-ordinary}.

\begin{figure}[!htbp]
    \centering
    \includegraphics[width=0.95\textwidth]{paperFigures/A2_two_panel_scaling_plus_geometry.pdf}
    \caption{Row-wise Lewis weight round counts on Gaussian and structured matrices (left), and leverage heterogeneity versus conditioning for the structured family (right).}
    \label{fig:structured-ordinary}
\end{figure}

Using the same $p$-grid, initialization, and stopping rule as the Gaussian experiment (\Cref{fig:A1}), the left panel of \Cref{fig:structured-ordinary}  shows that the tested structured family closely tracks the Gaussian baseline across the full tested range of $p$, requiring approximately $5$--$7\%$ more iterations for $p\geq7.3$. Thus the same qualitative $p$-dependence observed on the Gaussian baseline persists under this leverage-heterogeneous --- but well-conditioned --- departure from generic matrices.

\paragraph{Numerical robustness under ill-conditioning.}
We next test the floating-point implementation under increasing ill-conditioning while keeping the underlying Lewis map fixed. For each of five base instances $j\in[5]$, we construct $\bA_{j,s}=\bQ_j\bD_s\bV_j^\top$, where $\bQ_j\in\R^{100\times20}$ has orthonormal columns, $\bV_j\in\R^{20\times20}$ is orthogonal, and the diagonal entries of $\bD_s$ decrease geometrically; across nine levels, $\kappa_2(\bA_{j,s})$ ranges from approximately $1$ to $10^{16}$. Indeed, with $\bA_{j,0}=\bQ_j\bV_j^\top$,
\[
    \bA_{j,s}=\bA_{j,0}(\bV_j\bD_s\bV_j^\top),
    \qquad
    \calT_{\bA_{j,s}}(\bw)=\calT_{\bA_{j,0}}(\bw),
\]
since the Lewis map is invariant under invertible right transformations. Thus the intended exact family has the same trajectory at every conditioning level. 
The saved matrices are formed in floating-point arithmetic, so observed differences reflect end-to-end finite-precision effects, including both representation and map-evaluation error.

We measure this effect both at the level of one update and over the full trajectory. For fixed $p$, let $\bw_j^{(t)}(s)$ denote the $t^\mathrm{th}$ iterate for $\bA_{j,s}$, and let $T_j(0)$ denote the stopping time of the well-conditioned $s=0$ run. We record the following three quantities:
\[
    R_{j,t}(s):=\max_{i\in[m]}\left|
        \log\frac{w_{j,i}^{(t+1)}(s)}{w_{j,i}^{(t)}(s)}
    \right|,
    \quad
    E_{j,t}(s):=\max_{i\in[m]}\left|
        \log\frac{[\calT_{\bA_{j,s}}(\bw_j^{(t)}(0))]_i}
                  {[\calT_{\bA_{j,0}}(\bw_j^{(t)}(0))]_i}
    \right|,
    \quad
    E_{j,\max}(s):=\max_{0\le t<T_j(0)}E_{j,t}(s).
\]
Thus $R_{j,t}(s)$ measures progress along the run on $\bA_{j,s}$, and $E_{j,t}(s)$ compares the two floating-point maps at the reference iterate $\bw_j^{(t)}(0)$. Exact arithmetic gives $R_{j,t}(s)=R_{j,t}(0)$ and $E_{j,\max}(s)=0$.We reuse the primary $10^{-8}$ stopping tolerance from the row-wise scaling experiment (\Cref{fig:A1}) and test six values of $p$ spanning $4$ to $127.3$, for $270$ final trajectories. 
\begin{figure}[!htbp]
    \centering
    \includegraphics[width=0.95\textwidth]{paperFigures/A4_map_error_vs_conditioning.pdf}
    \caption{Finite-precision map discrepancy under increasing conditioning at representative values of $p$.}
    \label{fig:a4-map-error}
\end{figure}

As \Cref{fig:a4-map-error} shows, $E_{j,\max}(s)$ grows steadily with conditioning, reaching roughly $10^{-8}$ at $\kappa_2(\bA_{j,s})\approx10^8$. Still, the full iteration is quite stable through $s=27$ ($\kappa_2(\bA_{j,s})\approx1.3\times10^8$): all $30$ runs converge, with stopping times within about $2\%$ of their well-conditioned counterparts. At $s=33$ ($\kappa_2(\bA_{j,s})\approx8.6\times10^9$), none reaches $10^{-8}$ within $10^4$ updates. Thus, despite the growth of  $E_{j,\max}(s)$, the iteration count remains essentially unchanged over a broad conditioning range, until a finite-precision accuracy floor prevents the prescribed tolerance from being reached.

\begin{figure}[!htbp]
    \centering
    \includegraphics[width=.95\linewidth]{paperFigures/A4_posthoc_s33_residual_floor.pdf}
    \caption{Residual floors under ill-conditioning at representative values of $p$.}
    \label{fig:a4-residual-floor}
\end{figure}

\Cref{fig:a4-residual-floor} explains this failure. Across all $30$ runs, the residual falls below $10^{-6}$ and then fluctuates around a $10^{-7}$-scale floor; \Cref{fig:a4-residual-floor} shows three representative values of $p$. Thus finite precision imposes an accuracy floor above the stopping threshold, identifying a limitation outside the exact-arithmetic guarantee of \Cref{cor:ordinary-main}.

\subsubsection{Block Lewis weights: iteration counts}
\label{sec:expt-blkLW}

We finally ask whether the block Lewis weight iteration exhibits the same qualitative $p$-dependence as the row-wise iteration in \Cref{fig:A1}. We generate 50 matrices $\bA\in\mathbb{R}^{100\times20}$ with i.i.d.\ standard Gaussian entries, partition each into $k=20$ blocks of five rows, and reuse the same $p$-grid and primary $10^{-8}$ stopping rule as before. Keeping the blocks equal-sized and Gaussian ensures that we study the block iteration separate from heterogeneity in block size, rank, or conditioning.

\begin{figure}[t]
    \centering
    \includegraphics[width=.95\linewidth]{paperFigures/A3_main_body_style.pdf}
    \caption{Block Lewis weight round scaling with $p$ (left) and sensitivity to the stopping threshold (right). }
    \label{fig:block-scaling}
\end{figure}

The left panel of \Cref{fig:block-scaling} summarizes the $10^{-8}$ stopping time across the 50 matrices by its median and empirical 20th and 80th percentiles; the right reports the median at thresholds $10^{-4}$, $10^{-6}$, and $10^{-8}$. At all three thresholds, the iteration count increases steadily with $p$, while the spread across instances widens at large $p$. Thus the tested block instances provide empirical support for the $\widetilde{O}(p)$-type dependence proved in \Cref{thm:blkLW}. 


\subsection{Real-data experiments}\label{sec:real-data-expts}
We use our implementation of \Cref{alg:main} to run experiments on  the Intel Berkeley Research Lab sensor dataset~\cite{guestrin2004distributed} and meteorological data from the NOAA National Data Buoy Center. 

\subsubsection{Lewis weights and experimental design}
\label{sec:expts-design-geometry}
We now ask how block Lewis weights behave when used in experiment design. For the block Lewis weights $\bwbar(p)$, define the associated design allocation $
    \pi_i(p)
    :=
    \frac{\bar w_i(p)^{\,1-2/p}}
         {\sum_j \bar w_j(p)^{\,1-2/p}}.
$ 
These are precisely the normalized coefficients with which the block information matrices enter the regularized D-optimal-design formulation discussed in \Cref{sec:apps}. Recall that classical D-optimal design chooses an allocation to maximize the determinant of its information matrix; see, e.g., \cite{pukelsheim2006optimal}. The regularization favours less
concentrated allocations, and the resulting family interpolates from the
uniform design as $p\downarrow2$ toward scalar block D-optimal design as
$p\to\infty$. We first test whether this interpolation is visible on real
block geometries, before asking in \Cref{{sec:expts-block-failure}} how different
points along it behave under failures of different blocks. 

For each block $i\in[k]$, let $\bC_i\in\mathbb R^{r\times r}$ denote its positive-semidefinite contribution to the information matrix, so that a design $\bpi$ has information matrix $\sum_{i\in[k]}\pi_i\bC_i$. To quantify the above described interpolation, we track two complementary features of a design $\bpi$. Writing $L(\bpi):=\log\det\left(\sum_{i\in [k]} \pi_i \bC_i\right)$, define 
\[
    N_{\mathrm{eff}}(\bpi):=\frac{1}{\sum_{i\in[k]}\pi_i^2},
    \qquad
    G(\bpi):=\frac{L(\bpi)-L(U)}{L(D)-L(U)},
\]
where $U$ and $D$ denote the uniform and scalar block D-optimal designs, respectively. The effective support satisfies $N_{\mathrm{eff}}(U)=k$ and decreases as the allocation becomes more concentrated, while $G(U)=0$ and $G(D)=1$. Thus, in \Cref{fig:design-geometry}, moving left corresponds to concentrating the design, while moving up corresponds to recovering more of the available D-optimal log-determinant gain.

We evaluate this path on two real sensing datasets: the Intel Berkeley
Research Lab dataset \cite{guestrin2004distributed}, with one block per
sensor, and standard meteorological observations from the NOAA National
Data Buoy Center, with one block per Gulf buoy. These give
different block geometries on which to examine the same finite-$p$ design
path.

\begin{figure}[t]
    \centering
    \includegraphics[width=.95\linewidth]{paperFigures/B_design_geometry.pdf}
    \caption{Finite-$p$ block Lewis design paths on Intel and Gulf; square: uniform, star: D-optimal.}
    \label{fig:design-geometry}
\end{figure}

As \Cref{fig:design-geometry} shows, increasing $p$ traces a path from near-uniform toward near-D-optimal design: $N_{\mathrm{eff}}(\bpi(p))$ decreases across the evaluated $p$-grid, while $G(\bpi(p))$ increases. The interpolation occurs at different rates on the two datasets. For example, on Intel, $p=16$  attains $G=0.94$ with $N_{\mathrm{eff}}=30.2$, while the D-optimal endpoint has $N_{\mathrm{eff}}=23.5$; thus much of the available log-determinant gain is obtained before the allocation reaches the concentration of the D-optimal endpoint.

\subsubsection{Information preservation under block unavailability}
\label{sec:expts-block-failure}

The interpolation in \Cref{fig:design-geometry} suggests a natural
question: does the reduced concentration of finite-$p$ designs preserve
information when entire blocks become unavailable?  On the same Intel
sensor and Gulf buoy datasets considered above, we fix the design
allocation before availability is realized, and do not reallocate mass
from unavailable blocks.  For a surviving set $S$, write
\[
    M_S(\bpi):=\sum_{i\in S}\pi_i\bC_i,
    \qquad
    \Phi_S(\bpi):=\det(M_S(\bpi))^{1/r}.
\]
Our primary experiment conditions on exactly $f$ unavailable blocks and
chooses their identities uniformly, thereby separating failure severity
from which particular blocks happen to be unavailable.  We report the
mean surviving information
$\mu_f(\bpi):=\mathbb E[\Phi_S(\bpi)]$ together with the probability
$\rho_f(\bpi)$ that $M_S(\bpi)$ loses rank.

\begin{figure}[t]
    \centering
    \includegraphics[width=.95\linewidth]{paperFigures/C_exchangeable_mu_rho_combined_balanced.pdf}
\caption{Information preservation under block unavailability on Intel and Gulf datasets.}    \label{fig:block-failure}
\end{figure}

The top row of \Cref{fig:block-failure} shows the same qualitative
progression on both datasets.  When few blocks are unavailable, designs
near the concentrated, high-$p$ end retain the most information; as $f$
increases, the maximizing point moves steadily toward smaller
$p$.  Thus the amount of concentration favoured by the information
criterion changes with failure severity, consistent with the intuition that the
anti-concentration term acts as an  insurance
mechanism.

The bottom row of \Cref{fig:block-failure} separates this progression
from outright degeneracy.  Uniform and all finite-$p$ designs remain
full rank throughout the prespecified severity ranges, despite the
substantial crossings among their $\mu_f$ curves, so the migration
within the finite-$p$ family is not explained by rank loss.  The
D-optimal endpoint  begins to lose rank as failures become
severe, and its deterioration in that regime therefore combines reduced
surviving information with outright singularity.

Overall, the experiments indicate that different points along the
uniform-to-D-optimal Lewis path are favoured under different levels and
patterns of block unavailability.  We interpret this as descriptive
evidence consistent with the proposed anti-concentration mechanism on
these  datasets.

\section*{Acknowledgement}
We are extremely grateful to Guanghao Ye for numerous detailed  discussions about Lewis weights. We used ChatGPT interactively in developing and verifying our proofs, code, and write-up.  

\newpage 

\printbibliography

\newpage 
\appendix 
\section{Omitted proofs of supporting lemmas}\label{sec:app-technical-lemmas}
\blocklevscfact* 
\begin{proof}
The matrix $\bB(\bB^\top \bB)^{-1}\bB^\top$ is an orthogonal projector, and therefore $\bzero\preceq \bB(\bB^\top \bB)^{-1}\bB^\top \preceq \bI$. Restricting each matrix in this inequality chain to the principal submatrix of the rows corresponding to the $i^\mathrm{th}$ block in the statement of the lemma, we have $\bzero\preceq \bB_{[i]}(\bB^\top \bB)^{-1}\bB_{[i]}^\top\preceq \bI$. Therefore, the eigenvalues of $\bB_{[i]}(\bB^\top \bB)^{-1}\bB_{[i]}^\top$ all lie in $[0, 1]$. We therefore have \[\blevsc{i}(\bB)=\Tr\left( \bB_{[i]}(\bB^\top \bB)^{-1}\bB_{[i]}^\top\right)\leq \rank\left( \bB_{[i]}(\bB^\top \bB)^{-1}\bB_{[i]}^\top\right)\leq \rank(\bB_{[i]}).\]This proves the first assertion. For the second bound, we see that \[\sum_{i=1}^k \blevsc{i}(\bB) = \sum_{i=1}^k\Tr\left( \bB_{[i]}(\bB^\top \bB)^{-1}\bB_{[i]}^\top\right) = \Tr\left( \bB(\bB^\top \bB)^{-1}\bB^\top\right) = n.\]
\end{proof}

\begin{lemma}[Properties of $\bx\mapsto\log\det\bG(e^{\bx})$]\label{lem:logdetGexp-is-cvx} Following the notation of \Cref{sec:intro-notationPrelims}, let $\bA\in\R^{m\times n}$ be partitioned into nonzero row blocks $\bAblk{1}, \bAblk{2}, \dotsc, \bAblk{k}$, with $\bAblk{i}\in\R^{m_i \times n}$.  Given a vector $\bx\in\R^k$, define the function $f(\bx):= \log\det\bG(e^{\bx}) = \log\det\left(\bA^\top \Diag\left(e^{\omrp x_1}\bI_{m_1}, e^{\omrp x_2}\bI_{m_2}, \dotsc, e^{\omrp x_k}\bI_{m_k}\right)\bA\right)$. Then, the following properties hold. 
\begin{enumerate}
\item The function $f$ is convex. 
\item For every $i \in [k]$, we have that $\frac{\partial f(x)}{\partial x_i} = \omrp e^{\omrp x_i } \Tr\left(\bAblk{i}\bG(e^{\bx})^{-1}\bAblk{i}^\top\right) = \omrp \calT_i(e^{\bx}).$
\end{enumerate} 
\end{lemma}
\begin{proof}
When $\bA$ is a square matrix, the function $f$ is affine in $\bx$ and hence convex. For rectangular matrices, we instead invoke the Cauchy-Binet formula~\cite{horn2012matrix}. 

Given a matrix $\bB\in\R^{m\times n}$ with $m\geq n$, let $\bB_S\in\R^{n\times n}$ denote the square matrix obtained by selecting $n$ rows of $\bB$. Then the Cauchy-Binet formula says \[\det\left(\bB^\top\bB\right) = \sum_{\substack{S\subseteq[m]\\ |S| = n}} \left(\det(\bB_S)\right)^2.\numberthis\label{eq:logdetG-1}\] We proceed to apply this to $\bB : = \left(\Diag\left(e^{\omrp x_1}\bI_{m_1}, e^{\omrp x_2}\bI_{m_2}, \dotsc, e^{\omrp x_k}\bI_{m_k}\right)\right)^{1/2}\bA$. Consider a subset $S\subseteq [m]$ such that $|S|=n$. Let $z_i(S) := |S\cap \bI_i|$ denote those rows of block $\bI_i$ that are present in $S$. Then, using \Cref{eq:logdetG-1} in conjunction with the fact that the determinant of a diagonal matrix is simply the product of its entries, we have $\det\bG(e^{\bx}) = \sum_{\substack{S \subseteq[m]\\ |S| = n}} \det(\bA_S)^2 e^{\omrp \cdot \sum_{i = 1}^k x_i z_i(S)}$. Consequently, we may express $f$ as \[f(\bx) = \log \left(\sum_{\substack{S \subseteq[m]\\ |S| = n}} \det(\bA_S)^2 e^{\omrp \cdot \sum_{i = 1}^k x_i z_i(S)}\right).\] We now note that, as a function of $\bx\in\R^k$, this is a log-sum-exp of affine functions of $\bx$ and  is therefore convex~\cite[Section 3.1.5]{BV04}. For completeness, we now show the computation of the gradient, though this computation can be found, e.g., in \cite[Lemma 48]{lee2019solving}, \cite[Lemma 3]{fazel2022computing}. By the gradient of the determinant,  chain rule, and the cyclic permutation property of trace, we have \[\frac{\partial f(\bx)}{\partial x_i} = \Tr\left(\bG(e^{\bx})^{-1}\frac{\partial \bG(e^{\bx})}{\partial x_i}\right) = \Tr\left(\bG(e^{\bx})^{-1}\cdot \omrp e^{\omrp x_i} \bAblk{i}^\top \bAblk{i}\right) = \omrp e^{\omrp x_i}  \Tr\left(\bAblk{i} \bG(e^{\bx})^{-1} \bAblk{i}^\top\right) = \omrp \calT_i(e^{\bx}). \] 

\end{proof}

\section{Omitted proof on sensitivities}\label{sec:sens-proff}
\sensprop*
\begin{proof} The upper bound $\sens{i}(\bA)\leq n^{p/2-1}\wbar_i$ in the first part of the proposition is known \cite[Fact~3.3]{padmanabhan2023computing}. We therefore prove the reverse inequality $\wbar_i\leq\sens{i}(\bA)$.

Set $\bB\defeq\bWbar^{1/2-1/p}\bA$, and let $\bb_j^\top$ denote the $j^\mathrm{th}$ row of $\bB$. By the definition of Lewis weights, $\wbar_j=\levsc{j}(\bB)$ for every $j\in[m]$. We use the known fact that leverage scores are precisely $\ell_2$ sensitivities \cite[Section~2]{padmanabhan2023computing}, and hence
$\wbar_j=\max_{\bx\neq\bzero}\frac{|\bb_j^\top\bx|^2}{|\bB\bx|_2^2}$.

Fix $i\in[m]$, and let $\bx^{(i)}$ be a maximizer in the above expression for $\wbar_i$, scaled so that $\|\bB\bx^{(i)}\|_2=1$. Such a maximizer exists since $\bB$ has full column rank and the objective is homogeneous. By our choice of $\bx^{(i)}$, we have
$|\bb_i^\top\bx^{(i)}|^2=\wbar_i$.
Moreover, for every $j\in[m]$, since $\wbar_j$ is the maximum of the corresponding ratio and $\bx^{(i)}$ is one feasible choice with $\|\bB\bx^{(i)}\|_2=1$, we have
$|\bb_j^\top\bx^{(i)}|^2\leq\wbar_j$.

Define $z_j\defeq\frac{|\bb_j^\top\bx^{(i)}|^2}{\wbar_j}$ for every $j\in[m]$. The preceding inequalities imply $0\leq z_j\leq1$ for every $j$, while $|\bb_i^\top\bx^{(i)}|^2=\wbar_i$ implies $z_i=1$. Further,
$\sum_{j=1}^m\wbar_jz_j
=\sum_{j=1}^m|\bb_j^\top\bx^{(i)}|^2
=\|\bB\bx^{(i)}\|_2^2
=1$.

We now evaluate the $i^\mathrm{th}$ sensitivity at  $\bx^{(i)}$. Since $\bB=\bWbar^{1/2-1/p}\bA$, for every $j\in[m]$ we have $\bb_j^\top=\wbar_j^{1/2-1/p}\ba_j^\top$, or equivalently $\ba_j^\top=\wbar_j^{1/p-1/2}\bb_j^\top$. Therefore,
$|\ba_j^\top\bx^{(i)}|^p
=\wbar_j^{1-p/2}|\bb_j^\top\bx^{(i)}|^p$.
Using $|\bb_j^\top\bx^{(i)}|^2=\wbar_jz_j$, this becomes
$|\ba_j^\top\bx^{(i)}|^p
=\wbar_j^{1-p/2}(\wbar_jz_j)^{p/2}
=\wbar_jz_j^{p/2}$.
Consequently,

$$
\frac{|\ba_i^\top\bx^{(i)}|^p}{\|\bA\bx^{(i)}\|_p^p}
=
\frac{\wbar_i z_i^{p/2}}
{\sum_{j=1}^m\wbar_jz_j^{p/2}}.
$$

Since $p>2$ and $z_j\in[0,1]$, this implies $z_j^{p/2}\leq z_j$ for every $j$, and hence
$\sum_{j=1}^m\wbar_jz_j^{p/2}\leq\sum_{j=1}^m\wbar_jz_j=1$.
Since $z_i=1$, the preceding display therefore gives
$\frac{|\ba_i^\top\bx^{(i)}|^p}{\|\bA\bx^{(i)}\|_p^p}\geq\wbar_i$.
Finally, $\sens{i}(\bA)$ is the supremum of this ratio over all nonzero $\bx$, so
$\sens{i}(\bA)\geq\wbar_i$, concluding the proof. 
\end{proof}

\end{document}
